\scanpage{436}
\mspage{348}

Si donc \inserted{de plus} $\pi'>\pi\inserted{\,\geq\pi_0}$ (comme supposé au début), et si on suppose l'\deleted{\unsure{hyp.}}\unsure{HCG|HCC}, \deleted{\unsure{donc}} \unsure{donc|alors} la condition ${\pi'}^{\pi_0}=\pi'$ est automatique, et il reste la seule condition $\pi'\geq\pi''_1$ i.e. en définitive $\pi'\geq\mathrm{Card}\,\unsure{\mathrm{red}|\mathrm{rel}|\mathrm{réel}}\,\underline S_{\pi_0}$. Pour résumer :

\uline{Proposition} 4.13.27. Sous les conditions de 4.13.23 (\unsure{sans|sous} condition b) sur $\unsure{\Phi|\mathscr I|\mathscr D}$), soit $\underline S$ une sous-catégorie de $\unsure{M|\mathcal M}$ qui est $\Phi$-$\pi$-\unsure{adaptée|adapté}. Soit $\pi_0\leq\pi$ un cardinal \deleted{tel que} qui est utile pour $\underline S$ et pour $\underline M$, \deleted{\unsure{que ce}} et tel que $f:\underline C\to\underline M$ soit \unsure{$\pi_0$-accessible|$\pi_0$-accessibl} i.e. tel que $\underline S$ soit stable dans $M$ par $\varinjlim$ \unsure{grdes|grandes} devant $\pi_0$ (on peut p.ex. prendre $\pi_0=\pi$). Alors pour tout cardinal $\pi'>\pi$ tel que a) ${\pi'}^{\pi_0}=\pi'$ et b) $\pi'\geq\mathrm{Card}\,\mathrm{red}\,\underline S_{\pi_0}$), \deleted{\unsure{alors|on a}} $\underline S$ est $\Phi$-$\pi'$-accessible dans $M$.

Donc si $\pi'>\pi$ est tel que ${\pi'}^{\pi}=\pi'$, $\pi'>\mathrm{Card}\,\mathrm{red}\,\underline S_\pi$, alors $\underline S$ est $\Phi$-$\pi'$-\unsure{accessible|accessibl}

(Il suffit p.ex. de prendre $\pi'=2^c$, avec $c\geq\mathrm{Sup}\,(\pi,\mathrm{Card}\,\mathrm{red}\,\underline S(\pi))$, ou $\pi'=2^c$ avec $c\geq\pi$ et $2^c\geq\mathrm{card}\,\mathrm{red}\,\underline S(\pi)$.)

Si on admet l'\unsure{HCC|HCG}, il suffit donc qu'on ait $\pi'\geq\mathrm{Sup}\,(\pi,\mathrm{Card}\,\mathrm{red}\,\underline S(\pi))$.

% archive sheet number: 435

